\section{Physical robustness at logarithmic and polynomial thresholds}

\begin{definition}[A logarithmic truncation threshold]
    \label{def:aklv_logarithmic_threshold}
    \lean{TNLean.MPS.AreaLaw.aklvLogarithmicThreshold,
        TNLean.MPS.AreaLaw.aklvLogarithmicThreshold_le_div}
    \leanok
    For $g>0$, put
    \begin{align*}
        T(g)=1024+16\log(1/g).
    \end{align*}
    If $g\le1$, then $0\le T(g)\le1024/g$.
    Indeed, $\log(1/g)\le1/g-1$ and $1/g\ge1$.
\end{definition}

\begin{lemma}[Robustness above a logarithmic threshold]
    \label{thm:aklv_logarithmic_robustness}
    \lean{TNLean.MPS.AreaLaw.exists_aklv_logarithmic_robustness}
    \leanok
    \uses{def:aklv_normalization_regions,def:aklv_two_block_truncation,
        def:aklv_logarithmic_threshold}
    Let $s\ge6$, $a+s+1\le L$, and let $0\le h_i\le\Id$ act on
    $\{i,i+1\}$. Suppose the unit ground vector $\Omega$ of
    $H'=\sum_i h_i$ has energy $\mu$, actual first excited energy
    $\mu+g$, and
    $H'-\mu\Id\ge g(\Id-|\Omega\rangle\langle\Omega|)$, where $g>0$.
    Then $q\ge2$ and $g\le1$.
    The genuine reference minima can be chosen before $t$, with
    $E_0+\sum_j e_j=\mu$ and $0\le E_0\le6$.
    For every $t\ge T(g)$, the actual two-exterior truncation at
    $a+2,a+s-1$ has a phase-aligned unit ground vector $\phi$,
    gap at least $9g/10$, and
    \begin{align*}
        \|\phi-\Omega\|^2\le e^{-t/16}/100.
    \end{align*}
    This strengthens the sufficient threshold in
    \cite[Theorem~6.1]{Arad2013AreaLaw} within the stated window range.
\end{lemma}
\begin{proof}
    \leanok
    \uses{thm:aklv_two_block_scaled_low_cluster,
        thm:chain_positive_nearest_neighbor_gap_upper}
    The ground eigenvalue belongs to the spectrum because its unit
    eigenvector is nonzero. The two distinct spectral values force $q\ge2$,
    and the endpoint trial vector gives $g\le1$.
    Set $C_1=10^6\cdot64e^{3/4}$ and $\tau=e^{-t/16}$.
    The inequalities $\log C_1\le61$ and $t\ge1024-16\log g$ give
    \begin{align*}
        g+11 & \le t,                               &
        g+8\log\frac{C_1}{\tau g}
             & =\frac t2+g+8(\log C_1-\log g)\le t.
    \end{align*}
    Hence $0<\tau\le1$ and $t\ge\Theta(g,\tau)$.
    The actual spectral comparison supplies the required gap and squared
    phase error at most $\tau/100$. Multiplication by the inverse phase
    preserves the norm and eigenvalue and gives the asserted vector.
\end{proof}

\begin{lemma}[Robustness for a window with at least six interactions]
    \label{thm:aklv_physical_robustness}
    \lean{TNLean.MPS.AreaLaw.exists_aklv_physical_robustness}
    \leanok
    \uses{def:aklv_normalization_regions,def:aklv_two_block_truncation}
    Let $s\ge6$, $a+s+1\le L$, and let $0\le h_i\le\Id$ be supported on
    $\{i,i+1\}$. Suppose $\Omega$ is a unit ground vector of
    $H'=\sum_i h_i$ with energy $\mu$ and actual first excited energy
    $\mu+g$, where $g>0$, and suppose
    $H'-\mu\Id\ge g(\Id-|\Omega\rangle\langle\Omega|)$.
    The genuine reference minima $e_j$ can be chosen before $t$, with
    $E_0+\sum_j e_j=\mu$ and $0\le E_0\le6$.
    For every
    \begin{align}
        t\ge\frac{1024}{g}\left(\frac{E_0}{g}+1\right),
        \label{eq:aklv_physical_polynomial_threshold}
    \end{align}
    form the actual two-exterior truncation $B$ at cuts $a+2,a+s-1$
    with those same minima. It has an actual phase-aligned unit ground
    vector $\phi$, its gap is at least $9g/10$, and
    \begin{align}
        \|\phi-\Omega\|^2\le\frac{e^{-t/16}}{100}.
        \label{eq:aklv_physical_robustness_error}
    \end{align}
    This gives the polynomial-threshold and exponential conclusions of
    \cite[Theorem~6.1]{Arad2013AreaLaw} for the original bounded-term setting
    and the stated window range. The support and normalization corrections
    are recorded in \cite{gap:aklv13_two_block_robustness}.
\end{lemma}
\begin{proof}
    \leanok
    \uses{thm:aklv_logarithmic_robustness}
    Lemma~\ref{thm:aklv_logarithmic_robustness} gives $g\le1$ and
    the actual truncated ground vector above $T(g)$.
    Since $E_0\ge0$, the polynomial threshold implies
    $t\ge1024/g\ge T(g)$, so the conclusion follows.
\end{proof}

% The original nearest-neighbour/norm-one context is retained from source
% lines879--925. Theorem6.1 does not explicitly quantify over arbitrary independent
% exterior matrices. The present statement keeps s>=6 explicit; smaller windows
% remain separate, as documented by the cited source note.
